\providecommand{\R}{\mathbb{R}}
\providecommand{\Q}{\mathbb{Q}}
\providecommand{\N}{\mathbb{N}}
\providecommand{\Ccal}{\mathcal{C}}
\providecommand{\Fcal}{\mathcal{F}}
\providecommand{\Qcal}{\mathcal{Q}}

\section{Minimum diagrams and complete planar sector enumeration}
\label{pl:theory}

\subsection{The retained cone and its canonical lift}
\label{pl:canonical}

Fix one admissible quadrilateral sector.  Let \(t\) be the number of
tetrahedra and \(k\) the number of allowed quadrilateral coordinates.
The support reduction supplies a rational matching subspace
\(K\subseteq\R^k\), together with retained triangle classes partitioned
into vertex groups \(I_v\).  Write \(p_v=|I_v|\) and \(p=\sum_vp_v\).
The bound \(p\leq8k\) is a property of that reduction; the geometric
arguments below first keep \(p\) as a separate parameter.  Discarded
independent vertex-link factors can be reattached afterwards.  They add
their own link rays and do not change the nonlink rays considered here.

For each \(c\in I_v\), a rational linear potential
\(\pi_{v,c}:K\longrightarrow\R\) records the triangle-coordinate
differences imposed by matching.  In the actual construction the
potentials have integer coefficients, obtained by summing signed
quadrilateral labels along paths.  Most of the geometry only needs
rational coefficients.  Put
\[
 \Qcal=K\cap\R_{\geq0}^k,\qquad
 \Ccal=\left\{(q,x):
 \begin{array}{l}
 q\in\Qcal,\quad
 x_{v,c}=\pi_{v,c}(q)+a_v\quad(c\in I_v),\\
 a_v\in\R,\quad x_{v,c}\geq0
 \end{array}\right\}.
\]
The scalar \(a_v\) is shared within its group.  Let \(\ell_v\) be the
vector with zero quadrilateral part, entries one in group \(v\), and
zero in the other groups.  Define
\begin{equation}
 \mu_v(q)=\min_{c\in I_v}\pi_{v,c}(q),\qquad
 \Fcal(q)=
 \left(q,\bigl(\pi_{v,c}(q)-\mu_v(q)\bigr)_{v,c}\right).
 \label{pl:canonical-map}
\end{equation}
This is the canonical triangle normalization used in
quadrilateral-to-standard conversion: each vertex group has at least
one zero triangle coordinate.  Its role in normal-coordinate
conversion is established in Burton's work~\cite{Burton2009}.

\begin{lemma}[Canonical decomposition]
\label{pl:decomposition}
Every \(z=(q,x)\in\Ccal\) has a unique expression
\[
 z=\Fcal(q)+\sum_v b_v\ell_v,\qquad b_v\geq0.
\]
The map \(\Fcal\) is positively homogeneous and injective.
\end{lemma}
\begin{proof}
If \(x_{v,c}=\pi_{v,c}(q)+a_v\), nonnegativity is equivalent to
\(a_v+\mu_v(q)\geq0\).  Set \(b_v=a_v+\mu_v(q)\).
Subtraction gives the asserted expression.  The coefficients are
unique because they are the common differences within disjoint
groups; equivalently \(b_v=\min_{c\in I_v}x_{v,c}\).
The minimum of finitely many linear forms is positively homogeneous
for nonnegative scalars, proving homogeneity of \(\Fcal\).
Its first component is \(q\), proving injectivity.
\end{proof}

Set \(s(q)=\sum_{i=1}^kq_i\) and
\begin{equation}
 P=\{q\in\Qcal:s(q)=1\}.
 \label{pl:section}
\end{equation}
If \(\Qcal\neq\{0\}\), this is a nonempty compact rational polytope:
it is a closed subset of the standard simplex, and every nonzero
nonnegative \(q\) has \(s(q)>0\).  If
\(d=\dim K\leq3\), then the effective dimension
\(r=\dim P\) is at most two.  It may be smaller than \(d-1\).
An empty section, a point, and a segment must consequently be handled
as genuine possibilities even when the matching nullity is three.

Restrict each potential to \(\operatorname{aff}(P)\).  Identical
restrictions are identified for geometric purposes, without deleting
any output triangle coordinate.  For a group define its closed
minimum cells by
\begin{equation}
 C_{v,c}=
 P\cap\bigcap_{a\in I_v}
 \{\pi_{v,c}\leq\pi_{v,a}\}.
 \label{pl:min-cells}
\end{equation}
Take the full-dimensional cells, relative to
\(\operatorname{aff}(P)\), and all their faces.  Denote this minimum
subdivision by \(\mathcal D_v\), and denote their common refinement
by \(\mathcal D\).  In particular, the faces of \(P\) are present in
these complexes.  Let \(m_v\) be the number of full-dimensional
cells of \(\mathcal D_v\).

\begin{lemma}[Degenerate minimum labels]
\label{pl:degenerate-labels}
The full-dimensional minimum cells cover \(P\), meet face-to-face,
and satisfy \(1\leq m_v\leq p_v\).  A potential whose minimum cell
has smaller dimension introduces no additional geometric subdivision
of these cells and their faces.
\end{lemma}
\begin{proof}
Outside the finite collection of equality hyperplanes of distinct
restricted forms, the minimizing restricted form is unique.
Such points are dense in the relative interior of \(P\).
Approximating any point of \(P\) by them and taking a subsequence
with a fixed minimizing label proves coverage by closures of
full-dimensional cells.  Distinct such cells have disjoint relative
interiors.  Moreover
\[
 C_{v,c}\cap C_{v,a}
 =C_{v,c}\cap\{\pi_{v,a}-\pi_{v,c}=0\}.
\]
The difference on the right is nonnegative on \(C_{v,c}\);
thus the intersection is an exposed face of that cell, and
symmetrically of the other cell.

Suppose now that \(c\) is a label with a smaller-dimensional minimum
cell.  On any full-dimensional cell with minimizer \(a\), the affine
gap \(\pi_{v,c}-\pi_{v,a}\) is nonnegative.  Its zero set there is a
face.  More generally, on the relative interior of any face, an affine
nonnegative function can vanish only everywhere on that face or
nowhere in its relative interior.  Indeed, if it vanishes at a
relative interior point, its linear part must vanish in all directions
parallel to the face, since short displacements in both directions
remain in the face.  Hence the extra label cannot split an existing
face into smaller pieces.  The one-dimensional and zero-dimensional
versions of this argument are included.
\end{proof}

The common-cell extreme-ray characterization below is the general
minimum-envelope theorem of the preceding repository
report~\cite{MinimumEnvelopes}.  A proof is included to make explicit
the treatment of active quadrilateral coordinates and degenerate
minimum labels on which the present complete planar algorithm relies.

\begin{theorem}[Common-cell characterization]
\label{pl:ray-correspondence}
The extreme rays of \(\Ccal\) are its link rays
\(\R_{\geq0}\ell_v\) and the rays
\(\R_{\geq0}\Fcal(q)\) indexed by vertices \(q\) of \(\mathcal D\).
Different vertices give different nonlink rays.
\end{theorem}
\begin{proof}
The face with \(q=0\) is the nonnegative span of the link vectors.
Since their supports are disjoint, its extreme rays are precisely
the link rays.  It is a face of the whole cone: if nonnegative
quadrilateral parts sum to zero, both are zero.
For \(q\neq0\), any positive link coefficient in
Lemma~\ref{pl:decomposition} gives a decomposition into a nonzero
canonical vector and a nonzero vector with zero quadrilateral part.
These cannot be proportional.  A nonlink extreme ray must therefore
have its canonical representative.

Normalize \(s(q)=1\).  On each cell of \(\mathcal D\), a fixed
minimizing form may be chosen in every group, so \(\Fcal\) is affine
on that cell.  If \(q\) is not a vertex, it lies in the relative
interior of a positive-dimensional face.  There are distinct nearby
points \(q^+,q^-\) of that face with
\(q=(q^++q^-)/2\).  Therefore
\[
 \Fcal(q)=\tfrac12\Fcal(q^+)+\tfrac12\Fcal(q^-).
\]
The two vectors cannot be proportional, because their quadrilateral
parts are distinct and both have sum one.  Thus the ray is not
extreme.

Conversely, let \(q\) be a vertex of \(\mathcal D\).  Choose a
maximal common cell containing it and one defining minimizer \(c_v\)
for each group.  Suppose
\(\Fcal(q)=z_1+z_2\) with \(z_i\in\Ccal\).
The coordinate \(x_{v,c_v}\) is zero in \(\Fcal(q)\).
Nonnegativity forces the same coordinate to be zero in each
summand.  Thus \(c_v\) minimizes in both summands, both summands
are canonical, and their quadrilateral parts \(q_1,q_2\) belong
to the same chosen minimum cones.
A summand with \(q_i=0\) is zero: its zero coordinate in each group
forces every link coefficient to vanish.

If both summands are nonzero, set \(\lambda_i=s(q_i)>0\).
Then \(\lambda_1+\lambda_2=1\), and
\[
 q=\lambda_1(q_1/\lambda_1)+
   \lambda_2(q_2/\lambda_2).
\]
Both normalized points belong to the selected common cell.
A vertex of a polyhedral complex is a vertex of every cell containing
it, so both normalized points equal \(q\).
Positive homogeneity makes \(z_i=\lambda_i\Fcal(q)\), as required.
Working inside \(P\) has included all active inequalities \(q_i=0\);
they cannot be omitted from this extremality argument.
Finally, the normalization \(s(q)=1\) and the quadrilateral
projection distinguish the represented rays.
\end{proof}

\subsection{A planar count with the common boundary counted once}
\label{pl:counting}

Assume for this subsection that \(P\) is two-dimensional, and let
\(b\) be its number of genuine polygon corners.  Then \(b\leq k\).
We use polygon representations with no repeated vertex and no
collinear consecutive vertices.  In a subdivision graph, removable
collinear degree-two points are suppressed.  The following count is
useful because its bound for \emph{internal} edges is independent
of \(b\).

\begin{lemma}[One-group planar count]
\label{pl:euler}
For a minimum subdivision with \(m\) full-dimensional cells,
let \(V_i,V_b\) be its numbers of interior and boundary vertices,
and \(E_i\) its number of internal edges.  Then
\begin{equation}
 V_i+V_b-b\leq2(m-1),\qquad
 E_i\leq3(m-1).
 \label{pl:euler-bounds}
\end{equation}
The total number of polygon-vertex incidences over all its cells
is at most \(b+8(m-1)\).
\end{lemma}
\begin{proof}
Include the boundary cycle of \(P\) in the finite embedded graph,
and let \(c\geq1\) be its number of components.
There are \(m+1\) faces, including the exterior face, and
the number of boundary edges is \(V_b\).
Euler's formula gives
\[
 V_i+V_b-(E_i+V_b)+(m+1)=1+c,
 \qquad E_i=V_i+m-c.
\]
Every interior vertex has degree at least three.
There cannot be a degree-one interior endpoint in a face-to-face
cover by full-dimensional polygons.  If an interior vertex had
degree two, the two incident full-dimensional cells would share
their affine equality line through it; their two incident edges
would be collinear and the vertex removable.
Likewise a boundary vertex other than a genuine corner has degree
at least three, or it would be removable.
The \(b\) corners have degree at least two.  Summing degrees yields
\[
 2(E_i+V_b)\geq3V_i+3(V_b-b)+2b,
 \qquad
 2E_i\geq3V_i+V_b-b.
\]
Substituting the Euler identity proves
\[
 V_i+V_b\leq2m-2c+b.
\]
Since \(c\geq1\), this proves the first asserted bound.
Since \(V_b\geq b\), it also gives \(V_i\leq2m-2c\);
the Euler identity now gives \(E_i\leq3m-3c\).

An internal edge is counted twice in the boundaries of the cells,
and a boundary edge once.  Thus the total incidence count is
\(2E_i+V_b\), which is at most
\(6(m-1)+b+2(m-1)\).
This proof does not require a separate connectivity argument
for the subdivision graph.
\end{proof}

Put \(a_v=m_v-1\) and \(A=\sum_va_v\).
Let \(R\) denote the number of nonlink rays, equivalently the number
of vertices in the common subdivision.

\begin{theorem}[Planar output bound]
\label{pl:quadratic-bound}
If \(\dim P=2\), then
\begin{equation}
 R\leq b+2A+9\sum_{v<w}a_va_w.
 \label{pl:sharp-parameter-bound}
\end{equation}
The numbers of vertices, edges, and faces of the common subdivision
are all \(O(k+p^2)\).  For one nontrivial vertex group the sharper
bound \(R\leq b+2(m_v-1)\) holds.
\end{theorem}
\begin{proof}
The union of the polygon corners and the individual group vertices
contains at most \(b+2A\) points by Lemma~\ref{pl:euler}.
Every other overlay vertex is an intersection of internal edges
from different groups.  Two noncollinear segments meet in at most
one point.  Collinear overlapping segments produce no new vertex
except an endpoint of one of the original segments, already in
the preceding union.  Group \(v\) has at most \(3a_v\) internal
edges, so the additional vertex count is at most
\(\sum_{v<w}(3a_v)(3a_w)\).
Theorem~\ref{pl:ray-correspondence} transfers this count to rays.

There are at most \(3A\) internal segments before overlay.
Each can be split at only \(O(A)\) intersections and endpoints
of other segments, including collinear overlaps.
Consequently the overlay has \(O(A^2)\) internal edges.
Its boundary has \(b+O(A)\) edges, because each original internal
segment contributes at most two boundary endpoints.
Euler's formula bounds the number of faces by the same asymptotic
quantity.  If there is only one nontrivial group, the intersection
term in \eqref{pl:sharp-parameter-bound} vanishes.
\end{proof}

The distinction between internal edges and the common boundary is
essential in this proof.  Summing an \(O(k+p_v)\) bound for whole
group diagrams repeats the same boundary for each group and does
not by itself provide an \(O(k+p)\) collection of overlay segments.
Lemma~\ref{pl:euler} supplies the needed \(O(p_v)\) internal bound.

\begin{corollary}[Effective dimensions zero and one]
\label{pl:lowdim}
If \(P\) is a segment, then
\[
 R\leq2+\sum_v(m_v-1).
\]
If \(P\) is a point, \(R=1\); if \(P\) is empty, \(R=0\).
Thus every supplied sector of matching nullity at most three has
\(O(k+p^2)\) nonlink rays, and the retained normal-sector bound
\(p\leq8k\) makes this \(O(k^2)\).
\end{corollary}
\begin{proof}
On a segment, every minimum cell is an interval.
Distinct full-dimensional intervals have disjoint interiors,
and a group with \(m_v\) such intervals contributes at most
\(m_v-1\) interior breakpoints.
The common subdivision uses their union and the two section
endpoints.  The other cases follow directly from
Theorem~\ref{pl:ray-correspondence}.
\end{proof}

These minimum diagrams belong to established computational geometry.
For \(f_i(y)=u_i\cdot y+\beta_i\), put
\(z_i=-u_i/2\) and \(w_i=\|u_i\|^2/4-\beta_i\).  Then
\[
 \|y-z_i\|^2-w_i=\|y\|^2+f_i(y).
\]
The common quadratic term does not affect minimizers, so each group
diagram is a clipped power diagram~\cite{Aurenhammer1987}.
The present count is a direct planar argument adapted to the
normal-sector lift, rather than a new general theorem about
power diagrams.

\subsection{Exact construction and the cost of writing the rays}
\label{pl:construction}

One complete construction first obtains \(P\) in rational affine
coordinates.  If the section is two-dimensional, for each distinct
restricted group form \(f_c\) it computes
\[
 P\cap\bigcap_a\{f_a-f_c\geq0\}
\]
by successive exact half-plane clipping.
Empty or lower-dimensional results can be discarded immediately:
further intersections cannot restore their dimension.
Positive-area cells are recorded with their genuine corners.
Their edges on \(\partial P\) are omitted from the internal
segment set; the remaining edges are deduplicated.
The candidate vertices are the union of section corners, cell
corners, and intersections of internal edges from different groups.
All predicates are exact and all clipping inequalities are closed.
Multiple ties, parallel edges, coincident intersections, and
overlaps therefore require no generic perturbation.

For a segment, restrict the forms to a parameter \(u\in[0,1]\).
The lower envelope of lines can be constructed by sorting slopes
and maintaining increasing breakpoint positions.  Equal slopes
retain the lowest intercept, and equal affine forms retain one
representative.  Zero-length pieces are omitted, while the endpoints
of positive-length pieces are retained.  A point requires only one
canonical lift.  These branches use the feasible dimension, not
merely the dimension of the matching kernel.

\begin{proposition}[Constructive operation bound]
\label{pl:operations}
After the sector kernel has been constructed, complete geometry
and all native nonlink ray vectors can be obtained using
\[
 O(k^3+kp^2+p^3+(k+p)(k+p^2)+t(k+p^2))
\]
exact scalar arithmetic and comparison operations, with lower-order
terms absorbed.  Under \(p\leq8k\), this is
\begin{equation}
 O(k^3+t k^2).
 \label{pl:runtime}
\end{equation}
Integer gcd and lcm computations for primitive normalization are
additional operations on integers of polynomial bit length.
\end{proposition}
\begin{proof}
In an affine coordinate space of dimension at most two, one may
construct a bounded section by testing pairwise intersections of
the \(k\) boundary lines against all \(k\) inequalities and taking
their exact convex hull.  This elementary fallback costs \(O(k^3)\)
operations and handles empty, segment, and point sections as well.
The implementation can instead clip a justified coordinate box
when its chart uses actual \(q\)-entries.

A polygon initially having \(b\leq k\) sides has at most \(b+j\)
sides after \(j\) half-plane clips: a clip replaces a consecutive
boundary chain by at most one new side.
Thus computing every cell of group \(v\) takes
\(O(p_v^2(k+p_v))\) operations.
Summing gives \(O(kp^2+p^3)\).
There are only \(O(p)\) distinct internal segments, so their pairwise
intersection tests cost \(O(p^2)\).
Exact comparison sorting for deduplication stays within the displayed
bound.

Project all retained potentials into the affine chart once.
At each of \(O(k+p^2)\) resulting points, evaluating the quadrilateral
forms and triangle potentials and taking the group minima then
costs \(O(k+p)\), rather than a fresh length-\(k\) dot product for
every potential.  Expanding repeated triangle classes to the native
\(7t\)-coordinate vector, and checking its native matching equations,
costs \(O(t+k)\) per ray.  These charges prove the assertion.
\end{proof}

The quadratic ray count is not a quadratic bound for explicit
native output.  The geometry uses \(O(k+p^2)\) candidate points,
but writing \(R\) native vectors requires \(\Theta(tR)\) coordinate
positions in this representation.  In the normal-sector regime,
the output alone can therefore have \(O(tk^2)\) entries.
This distinction is retained in \eqref{pl:runtime}.

In the integer-potential model, one may first clear the denominators
of the quadrilateral direction and divide its entries by their gcd.
Every triangle potential, and then every canonical triangle
coordinate, is integral on this primitive integer \(q\).
The whole standard vector is primitive because its quadrilateral
coordinates already have gcd one.
This procedure uses \(O(k)\) gcd/lcm steps per ray, besides scalar
evaluation and output.  For arbitrary rational potentials the
corresponding safe procedure clears denominators in the entire
lift instead; the ray theorem does not assume that primitive
quadrilaterals imply integral triangles.

Let \(B\) bound the bit length of the rational affine coefficients
after exact coordinate reduction.  Every geometric vertex is the
intersection of two source boundary or active minimum-equality
lines.  The constant-size determinant formulas therefore have
polynomial, indeed \(O(B)\), rational encoding length up to fixed
constant factors.  Section endpoints obey the analogous
one-dimensional bound.  The corresponding bound for exact linear
algebra does not assume that numerical coefficients remain small
as \(t\) grows.  After clearing row denominators, suppose the integer
matrix entries have at most \(\beta\) bits in absolute value.
Hadamard's inequality bounds every rank-\(\rho\) minor by
\[
 |\det M|\leq(\sqrt{\rho}\,2^\beta)^\rho,\qquad
 \log_2|\det M|\leq\rho\beta+
                  \tfrac{\rho}{2}\log_2\rho
\]
when the determinant is nonzero.
Exact kernel bases are obtained from ratios of such minors.
The normalized chart then uses constant-size inversions because
\(d\leq3\); projecting a source potential adds exact products and
sums over its \(k\) coefficients.
These estimates charge both matrix dimensions and coefficient bit
lengths, including row denominators.  In the generic rational model,
these are part of the supplied data.
For native matching input, path potentials sum \(O(t)\) incidence
labels, so their coefficient bit lengths are polynomially bounded
before applying the determinant estimate.
These determinant bounds consequently give polynomial bit length
uniformly for growing \(t\) and \(k\), including families with
exponentially large primitive coordinates.
Retaining support-line descriptions also makes the geometric
bounds transparent during clipping.  Clearing denominators of \(k\)
coordinates can produce a common denominator with \(O(kB)\) bits
in the generic rational model.  This remains polynomial, but its
arithmetic is not a unit-cost bit operation.
Thus the construction has polynomial bit complexity for this fixed
dimension, with rational arithmetic and integer normalization costs
charged at their actual operand lengths.

\begin{remark}[Inactive equalities are not minimum edges]
\label{pl:inactive}
It is not correct to treat every equality of two potentials as
an edge of the canonical lift.
On \(P=\{(x,y):x,y\geq0,\ x+y\leq1\}\), consider
\(0,1+x,1+y\).  The latter two forms agree on \(x=y\), but
neither is minimal anywhere.  The canonical lift is affine on
the whole triangle.  The intersection \((1/2,1/2)\) of that
irrelevant equality with the boundary is not an extreme lifted
ray; it is a convex combination of the adjacent corner lifts.
An arrangement of all pairwise equalities may be used only with
a separate extremality filter, and it loses the active-skeleton
output bound proved above.
\end{remark}

\section{Independent certificates of complete minimum coverage}
\label{cv:coverage}

The producer's polygon-clipping sequence need not belong to the
verification trust boundary.  A compact alternative submits the
resulting minimum cells and proves their coverage by exact area
or interval length.  The checker reconstructs the matching
subspace and all native corner potentials independently, so the
following statements concern source functions, not functions
accepted from the producer.

\subsection{Certifying the affine chart and the section}
\label{cv:chart}

\begin{proposition}[Chart completeness]
\label{cv:chart-completeness}
Suppose the source checker establishes \(\dim K=d\), where
\(1\leq d\leq3\).  A submitted affine chart
\[
 q(y)=c_0+\sum_{j=1}^{d-1}y_jc_j
\]
parametrizes all of \(K\cap\{s=1\}\) if the checker verifies
\(c_j\in K\), \(s(c_0)=1\), \(s(c_j)=0\) for \(j>0\),
and distinct coordinate indices \(i_1,\ldots,i_{d-1}\) satisfying
\[
 q_{i_j}(y)=y_j.
\]
Every feasible chart point then lies in \([0,1]^{d-1}\).
\end{proposition}
\begin{proof}
The linear map
\[
 q\longmapsto(s(q),q_{i_1},\ldots,q_{i_{d-1}})
\]
sends \(c_0,c_1,\ldots,c_{d-1}\) to the standard basis of
\(\R^d\).  These vectors are linearly independent.
Since they belong to the \(d\)-dimensional source kernel, they
form its basis.  A vector in that kernel with \(s(q)=1\)
therefore has the asserted affine expression, with coefficients
\(q_{i_j}\).  For a nonnegative vector of coordinate sum one,
each of those actual coordinates belongs to \([0,1]\).
\end{proof}

If \(d=0\), or if \(s\) vanishes identically on \(K\), the
nonnegative cone contains only zero and no chart is needed.
Otherwise an affine chart is required even when its feasible
section later turns out to be empty.  Nullity alone cannot
justify accepting a missing section.

\begin{proposition}[Source-facet section certificate]
\label{cv:section-facets}
Let a two-dimensional chart describe the source section by
affine inequalities \(h_i(y)\geq0\).
A submitted nondegenerate convex polygon \(S\), given by its
genuine corners in cyclic order, equals the source section \(P\)
if the following checks pass:
every vertex satisfies all source inequalities, and every edge
of \(S\) lies on the zero line of a nonconstant source form \(h_i\).
\end{proposition}
\begin{proof}
Vertex feasibility and affinity give \(S\subseteq P\).
For an edge of \(S\), its checked source form vanishes on the
edge line and is nonnegative on \(S\).
The form is nonconstant and \(S\) has interior, so it is positive
at some point of \(S\).  Consequently its nonnegative half-plane
is exactly the inward edge half-plane of \(S\).
Every point of \(P\) lies in that half-plane.
Intersecting the inward half-planes for all edges gives
\(P\subseteq S\), proving equality.
\end{proof}

This two-sided argument does not reconstruct a clipping history.
If the submitted section has fewer than three vertices in a
two-dimensional chart, the checker instead enumerates all feasible
pairwise source-line intersections within the established coordinate
box and takes their hull.  A nonempty bounded polyhedron in the plane,
including a segment or a point, has extreme points determined by
independent active constraints, so the fallback is complete.
In a one-dimensional chart, direct exact lower and upper endpoint
bounds suffice.

\subsection{The area and length certificate}
\label{cv:measure}

\begin{theorem}[Exact minimum-cell coverage certificate]
\label{cv:area-theorem}
Let \(P\) be a nondegenerate convex polygon and let
\(f_1,\ldots,f_n\) be its source affine functions after identifying
identical restrictions.  Suppose a certificate submits finitely
many positive-area convex polygons \(S_i\), each carrying one
source label, subject to the following conditions:
\begin{enumerate}
 \item no restricted affine function labels two submitted polygons;
 \item every submitted vertex belongs to \(P\) and satisfies
       \(f_i\leq f_j\) for every source function \(f_j\);
 \item the sum of submitted polygon areas equals
       \(\operatorname{area}(P)\).
\end{enumerate}
Then the polygons are exactly the full-dimensional source
minimum cells.  In particular they cover all of \(P\), including
its boundary, and no full-dimensional minimum cell is omitted.
\end{theorem}
\begin{proof}
The vertex checks extend to each polygon by affinity and convexity,
so \(S_i\subseteq C_i\), where \(C_i\) is the true minimum cell.
If the relative interiors of polygons carrying different labels
intersected, they would intersect in a nonempty open set.
On that set both \(f_i\leq f_j\) and \(f_j\leq f_i\) would hold.
The affine difference would vanish on an open set and hence
identically, contrary to the distinct-label condition.
Thus pairwise intersections have area zero, and
\[
 \operatorname{area}\left(\bigcup_iS_i\right)
 =\sum_i\operatorname{area}(S_i)
 =\operatorname{area}(P).
\]

The union is finite and closed.  If a point \(z\in P\) were
missing, some neighborhood of \(z\) would be disjoint from the
union.  Its intersection with \(P\) has positive area, even
when \(z\) is on the boundary.  To see this last statement, choose
an interior point \(w\in P\).  Points
\((1-\varepsilon)z+\varepsilon w\), for sufficiently small
\(\varepsilon>0\), lie in that neighborhood and in the interior
of \(P\); a sufficiently small disk around one such point lies
in both.  A missing point therefore contradicts area equality.
The polygons cover all of \(P\), not merely an almost-everywhere
subset.

Now let \(C_i\) be a true full-dimensional minimum cell.
Its uniquely minimized points are dense in it: within its
interior, remove the finite union of equality lines with the other
distinct source forms, and then take the closure.
At such a unique-minimizer point, coverage is possible only by a
polygon carrying label \(i\).
This proves that label \(i\) is submitted, and closedness gives
\(C_i\subseteq S_i\).
Together with the already established reverse inclusion it yields
\(C_i=S_i\).
Conversely every submitted polygon has positive area and is
contained in its source minimum cell, so that source cell is
full-dimensional.  These two conclusions prove exact equality of
the cell inventories.
\end{proof}

\begin{corollary}[Interval coverage certificate]
\label{cv:length-theorem}
Parametrize a nondegenerate source segment by \(u\in[0,1]\).
The corresponding statement holds for positive-length submitted
intervals, with pairwise distinct restricted affine labels,
endpoint dominance, and total parameter length one.
\end{corollary}
\begin{proof}
Endpoint dominance holds throughout an interval.
Distinct affine functions on a line cannot agree on a nonempty
open interval, so submitted interiors are disjoint.
Exact total length proves that their finite closed union is
\([0,1]\); any missing point would leave a positive-length relative
neighborhood.  Density of unique minimizers then proves equality
with the true full-dimensional minimum intervals, as above.
\end{proof}

For use as a vertex certificate, a submitted polygon is additionally
required to equal the canonical convex hull of its listed vertices,
with positive signed area.  This verifies its cyclic order and
removes repeated or collinear entries.  A test of consecutive
turns alone is insufficient, since a self-intersecting polygon
can have consistently signed local turns.
The one-polygon-per-function requirement is also essential:
without it, a genuine minimum cell could be divided into artificial
pieces while preserving both dominance and total area.
Theorem~\ref{cv:area-theorem}, together with the canonical hull
check, ensures that the submitted vertices are precisely genuine
cell corners.

Groups with only one distinct restricted function have the whole
section as their minimum cell and need no separate geometry record.
For the remaining source groups, the checker verifies the inventory
of group identities.  It permits omission of lower-dimensional
minimum labels, as justified by Lemma~\ref{pl:degenerate-labels}.
On a feasible segment inside a two-dimensional chart, equality of
restricted functions is tested by their values at the two endpoints;
equality of their ambient coefficient tuples would be too strong.

\begin{theorem}[Complete ray-list replay]
\label{cv:complete-replay}
After independently validating the source model, chart, and section,
suppose the checker validates every nontrivial group's coverage
certificate as above.  It forms the union of section and cell
corners, adds intersections of internal edges from different
groups, and canonically lifts every resulting point using the
independently reconstructed source potentials.
Exact equality with the submitted ray list certifies that list as
all nonlink standard rays of the supplied sector.
\end{theorem}
\begin{proof}
The coverage theorems identify the submitted polygons or intervals
with the exact source minimum cells, including all cells with
nonempty relative interior.  Their genuine corners and internal
edges consequently form the exact individual group skeletons.
The overlay argument in Theorem~\ref{pl:quadratic-bound} identifies
the constructed point set with all vertices of their common
refinement; collinear overlaps need only endpoints already present.
The lower-dimensional cases are covered by
Corollary~\ref{pl:lowdim}.
Finally Theorem~\ref{pl:ray-correspondence} identifies the
independently lifted vectors with exactly the required rays.
\end{proof}

\subsection{Certificate size and verification complexity}
\label{cv:cost}

If \(g\) groups require geometry records, Lemma~\ref{pl:euler}
bounds the number of polygon-vertex incidences by
\[
 \sum_v\bigl(b+8(m_v-1)\bigr)=gb+8A.
\]
Thus the geometry certificate has \(O(k^2)\) vertices when
\(p\leq8k\); this charge includes repeated copies of the section
boundary.  The complete certificate also includes the native
ray list, with \(O(tk^2)\) scalar entries.
Both statements count entries, not their integer bit lengths.

After independent source reconstruction and rank computation,
verification of a correct certificate takes \(O(k^3+tk^2)\)
rational arithmetic and comparison operations.
Indeed, feasibility and source-boundary scans cost
\(k\sum_vO(b+m_v)=O(k^3)\).
If \(t_v\) is the number of native triangle corners in group \(v\),
dominance tests cost
\[
 \sum_v O((b+m_v)t_v)=O(tk),
\]
because \(\sum_vt_v=4t\) and \(b,m_v=O(k)\).
Canonical convex-hull checks and segment intersections fit within
the cubic bound, and reconstructing \(O(k^2)\) native rays costs
\(O(tk^2)\), with integer normalization separately charged as
above.  The cost of reconstructing the native model and proving
its rank is additional, not hidden in this geometric replay bound.

For arbitrary malformed certificates, their supplied lengths and
integer bit lengths must also be charged.  In particular an
unrestricted submitted vertex list can require work to parse and
reject before its geometric size is known.  The displayed
parameter bound describes correct certificates, or certificates
subject to suitable early size bounds; it does not make parsing
an arbitrarily long invalid input free.

\section{A sharp abstract quadratic family}
\label{cf:grid}

The quadratic dependence on the number of potentials is already
necessary in a three-dimensional quadrilateral cone.
The following family belongs to the algebraic cone-and-potential
model.  No realization by the matching equations of a triangulated
three-manifold is asserted.

\begin{proposition}[Exact grid count]
\label{cf:grid-count}
For every integer \(m\geq2\), there is a rational model with
\(\dim K=3\), two groups of \(m\) potentials each, and exactly
\[
 p=2m,\qquad R=m^2+2m
\]
nonlink extreme rays.  Its potential coefficients are integers
of \(O(\log m)\) bits.
\end{proposition}
\begin{proof}
Take \(K=\R^3\), \(\Qcal=\R_{\geq0}^3\), and write
\(q=(x,y,1-x-y)\) on its normalized triangle
\[
 P=\{(x,y):x\geq0,\ y\geq0,\ x+y\leq1\}.
\]
For \(j=0,\ldots,m-1\), define
\[
 \pi^A_j(q)=j^2s(q)-8mj\,q_1,\qquad
 \pi^B_j(q)=j^2s(q)-8mj\,q_2.
\]
On \(s(q)=1\), the first group consists of the functions
\(j^2-8mjx\).  Completing the square gives
\[
 j^2-8mjx=(j-4mx)^2-16m^2x^2.
\]
The minimizing index is therefore the nearest allowed integer to
\(4mx\), with ties between adjacent indices at
\[
 \alpha_j=\frac{2j+1}{8m},
 \qquad j=0,\ldots,m-2.
\]
All \(m\) minimum cells have nonempty interior.
The internal edges are precisely the \(m-1\) vertical segments
\(x=\alpha_j\) clipped to \(P\).
The second group gives the \(m-1\) corresponding horizontal
segments \(y=\alpha_j\).

Every pair of one vertical and one horizontal segment meets in
the interior of \(P\), since
\(\alpha_i,\alpha_j<1/4\), so
\(\alpha_i+\alpha_j<1/2\).
These give \((m-1)^2\) distinct interior vertices.
Each vertical segment has two boundary endpoints, and so does
each horizontal segment, giving \(4(m-1)\) boundary vertices.
They are distinct from the three triangle corners.
Nor can a vertical and horizontal endpoint coincide on the
hypotenuse: that would require
\(\alpha_i+\alpha_j=1\).
There are no other group vertices or overlay crossings.
Hence the complete vertex count is
\[
 3+4(m-1)+(m-1)^2=m^2+2m.
\]
Theorem~\ref{pl:ray-correspondence} gives exactly that many
nonlink extreme rays.  The largest potential coefficient has
absolute value \(O(m^2)\), proving the bit-length assertion.
\end{proof}

This proves sharp quadratic order for the abstract \(O(k+p^2)\)
bound, not optimality of the constants in
\eqref{pl:sharp-parameter-bound}.
If desired, the model can be padded with identically zero
quadrilateral coordinates to width \(k=\max(3,m)\); its effective
cone and ray count remain unchanged and \(p\leq8k\) then holds.
Such padding still establishes no topological realization.
Determining which quadratic examples occur in genuine normal
matching systems is a separate question.  Likewise, a complete
polynomial solver for each supplied low-nullity sector gives no
bound here on how many sectors a global recognition procedure
must inspect.
